\section{Data and provenance audit}\label{sec:data}

\paragraph{Sources.}
Model development used an archive of AGMARKNET daily reports for 2001-01-10 to 2026-04-21: a Kaggle dataset (\texttt{khandelwalmanas/daily-commodity-prices-india}, version 134, GODL-India).
It has \num{75984017} rows; after quality classification and business-key de-duplication, \num{75644801} rows remain.
Its columns match the official data.gov.in resource \emph{Variety-wise Daily Market Prices Data of Commodity}.
After development, an API key became available. We back-filled \num{7720662} official records for 2025-10-01 to 2026-10-22, one request series per date, 387 dates with no failure.
Our second layer combines the archive before 2025-10-01 with the official record afterwards, for \num{82114083} rows in total.

\paragraph{Scope and panel.}
We use 36 staple commodities: vegetables, fruits, cereals, pulses, oilseeds and cotton.
They were chosen from training-period data:
\begin{enumerate}
  \item take an a-priori list of crop-name patterns;
  \item keep matches with more than \num{50000} valid reports in 2021--2023;
  \item drop two duplicate or processed forms;
  \item keep the 36 most-reported.
\end{enumerate}
An entity (series) is (state, district, market, commodity, variety).
It is in the panel if it has at least 300 report days in 2021--2023, which gives \num{11845} series and \num{21.06}M feature rows.
A causal sanity rule removes reports more than a factor of five from the geometric mean of the previous ten reports (\num{28567} reports, 0.086\%).

\paragraph{Identifier drift.}
Several changes in how identifiers are published would break series if left untreated:
\begin{itemize}
  \item many series switch their grade label from \emph{FAQ} to \emph{Local} in January 2025. No (state, district, market, commodity, variety, date) key carries two grades in the archive since 2023, so grade is an attribute. The official record is different: there, 1.2\% of entity-days (\num{85315}) carry several grades, and a deterministic rule keeps one of them;
  \item markets gain an `` APMC'' suffix with the AGMARKNET~2.0 portal (live 2025-11-07);
  \item the official API writes no space before ``('' and uses ``Keralam'' for Kerala.
\end{itemize}
One normalisation rule, shared by the batch and streaming paths, maps all of these. On 2026-04-20, all \num{5152} valid archive keys are found among the \num{22716} official records.

\paragraph{Mirror completeness.}
Table~\ref{tab:mirror} compares daily volumes.
On the sample days in Table~\ref{tab:mirror} the mirror holds 94--100\% of official reports before October 2025 and 0--54\% afterwards. By month, it holds 70.4\% in October 2025 and 5.9--38.9\% from November 2025 to April 2026, which is 30.0\% overall for 2025-10-01 to 2026-04-21.
Every mirror record in that window exists in the official data, and 99.33\% have an identical modal price. The mirror is incomplete, not corrupt.
Missing reports lengthen the gaps between consecutive reports of a series:
\begin{itemize}
  \item the median gap is 3 days in the mirror and 1 day in the official record;
  \item 16.2\% of mirror gaps exceed a week, against 2.4\% in the official record.
\end{itemize}
Longer gaps push ``next report'' targets further ahead and thin the trailing windows the features rely on.

\begin{table}[t]\centering\small
\caption{Daily report counts, official API versus the public mirror (sample days).}\label{tab:mirror}
\begin{tabular}{@{}lrrr@{}}\toprule
Date & Official & Mirror & Share\\\midrule
2025-03-12 & 22,040 & 21,967 & 0.997\\
2025-09-17 & 20,703 & 19,542 & 0.944\\
2025-10-15 & 21,899 & 20,611 & 0.941\\
2025-11-20 & 17,824 & 0 & 0.000\\
2025-12-17 & 19,566 & 10,559 & 0.540\\
2026-02-18 & 18,979 & 5,992 & 0.316\\
2026-04-20 & 22,716 & 5,162 & 0.227\\\bottomrule
\end{tabular}\end{table}
